\documentclass[10pt,a4paper]{amsart}
\usepackage[margin=19mm]{geometry}
\usepackage{amsmath,amssymb,mathtools}
\usepackage[scr=rsfs,cal=euler]{mathalfa}
\usepackage{xfrac}
\usepackage[unicode,hidelinks,bookmarks=false]{hyperref}
\newcommand{\ie}{i.e.\ }
\newcommand{\cf}{cf.\ }
\newcommand{\resp}{respectively\ }
\newcommand{\Subsection}{\subsection}
\providecommand{\nobreakdash}{\nobreak\hbox{-}}
\setlength{\emergencystretch}{2em}
\multlinegap0pt
\usepackage{amssymb}
\usepackage{xcolor}
\newtheorem{theorem}{Theorem}
\newtheorem{corollary}{Corollary}
\def\P{{\mathcal P}}
\def \d{{\delta}}
\def \e{{\varepsilon}}
\def \k{{\kappa}}
\def \m{{\mu}}
\def \o{{\omega}}
\def \p{{\varphi}}
\def \qq{{\qquad}}
\title{Divisibility and primality in random walks}
\author{Michel J.\,G.\, WEBER}
\date{Source: arXiv:2310.14828v2 (CC BY 4.0)}
\begin{document}
\maketitle

\noindent\emph{Source and licence.} Divisibility and primality in random walks. Version arXiv:2310.14828v2 (CC BY 4.0), distributed by arXiv under the Creative Commons Attribution 4.0 International licence, http://creativecommons.org/licenses/by/4.0/, as recorded in the arXiv licence metadata for this exact version and shown on the abstract page https://arxiv.org/abs/2310.14828v2. Full licence notice: \texttt{licenses/arxiv-2310.14828-CC-BY-4.0.txt}.\par\smallskip

\noindent\textit{Editorial scope.} Counted unit: the article's theorem bnmbound (source0.tex lines 1940--1948). The article's complete original proof of it is delivered with it (source0.tex lines 1949--2057, one \texttt{proof} environment of 109 lines). No mathematics is written, repaired or re-derived: the statement and the proof are the article's own lines, in the article's own notation, and no retained line is substituted or re-flowed.  Retained uncounted as the source-local context the proof needs: lines 1735--1744; lines 1923--1931.  Nothing was omitted from any retained span, so the package carries no omission record.  No retained line contains a citation, so the wrapper carries no bibliography and no bibliography entry.  No retained line contains a figure environment, an image-inclusion command or drawing code, so no image is delivered and the package declares no asset dependency.  Editorial formatting only, in addition: an AMS \texttt{amsart} preamble that restates the article's own declarations and macros used by the retained lines, namely the environments \texttt{theorem}, \texttt{corollary} on separate counters, together with the standard packages those lines need.  The retained environments print in this document as Theorem 1, Corollary 1 in order of appearance, whereas the article prints the same items with the numbering of its own full document.  The complete native source of the article as distributed in the arXiv e-print for this exact version is bundled unchanged as \texttt{sources/148924/source0.tex}.
 \begin{theorem}\label{p4.5} For any $\e>0$, there exists a constant $C_\e$ depending on $\e$ only, such that  for any
 positive  integers
 $n,m, D$ with 
 $m>n$
\begin{equation}\label{4.11} \Big|\P \{ D|B_nB_m \}-  {1\over D  2^{ m-n}} \sum_{k=0}^{m-n} { {{m-n}\choose {k}}     } \varrho_k(D) 
\Big|
\le C_\e
\left({     D  ^{1 + \e}\over   n }\right)^{1/2 }. 
\end{equation} 
\end{theorem}

\begin{corollary}\label{rokest} For any positive integers  $D,k$,   
$$\varrho_k(D)\le 2^{\o ((k, D_{1/2}))}\, (k, D_{1/2})
%2^{\o(D) }(k\wedge  \sqrt D) 
 \qq
  {\rm and}\qq
\varrho_0(D)\le  
\sqrt D.
$$   
\end{corollary}  

\begin{theorem} \label{bnmbound} We have for any positive integers $n,m,D$,  
\begin{eqnarray*}   \P \{ D|B_nB_m \} 
 &\le &{1  \over D }\sum_{d| D_{1/2} } 2^{\o (d)}d
 \Big(\sum_{u|  (D_{1/2}/d )}\m(u)\, \P\{ud|B_{m-n}\}\Big)
\cr &  &   +{1\over 2^{  m-n }\sqrt D } +C_\e
\left({     D  ^{1 + \e}\over   n }\right)^{1/2 }. 
\end{eqnarray*} 
%$$  \P \{ D|B_nB_m \}    \le   {1\over 2^{  m-n }\sqrt D}+    {2^{\o(D)}( m-n)  \over D } + C_\e {     D  ^{{1/ 2} + \e}\over   \sqrt n }  .    $$
   \end{theorem}  

\begin{proof} By Theorem \ref{p4.5},
 for any $\e>0$, there exists a constant $C_\e$ depending on $\e$ only, such that  for any
 positive  integers
 $n,m, D$ with 
 $m>n$
\begin{equation}\label{4.11a} \Big|\P \{ D|B_nB_m \}-  {1\over D  2^{ m-n}} \sum_{k=0}^{m-n} { {{m-n}\choose {k}}     } \varrho_k(D) 
\Big|
\le C_\e
\left({     D  ^{1 + \e}\over   n }\right)^{1/2 }. 
\end{equation} 
 By Corollary \ref{rokest}, 
 %and since 
 %$ (k, D_{1/2})= \sum_{d|k\,|\,d|D_{1/2}}\p(d)$,
\begin{eqnarray}\label{323}
& & {1\over D  2^{ m-n}}  \sum_{k=0}^{m-n}   {{m-n}\choose {k}} \,  \varrho_k(D)
\ = \  {\varrho_0(D)\over 2^{ 
m-n }D }+{1\over D  2^{ m-n}} 
\sum_{k=1}^{m-n}   {{m-n}\choose {k}}\,
 \varrho_k(D)  
\cr &  \le  &  {1\over 2^{  m-n }\sqrt D }+  {1  \over D 2^{  m-n }}  \sum_{k=1}^{m-n}   {{m-n}\choose {k}} 2^{\o ((k, D_{1/2}))}\, (k, D_{1/2})
\cr & =  &  {1\over 2^{  m-n }\sqrt D }+ 
 {1  \over D 2^{  m-n }}\sum_{d| D_{1/2}} 2^{\o (d)}d
 \sum_{k=1\atop (k, D_{1/2})=d}^{m-n}   {{m-n}\choose {k}}  
%\cr &  =  &  {1\over 2^{  m-n }\sqrt D }+  {1  \over D 2^{  m-n }} \sum_{ d|D_{1/2}} \, \p(d)\sum_{k=1\atop d|k}^{m-n}   C^k_{m-n} 2^{\o ((k, D_{1/2}))}%\cr &  \le  &  {1\over 2^{  m-n }\sqrt D }+  {2^{\o(D)}  \over D 2^{  m-n }}  \sum_{k=1}^{m-n}   C^k_{m-n}k
 %\cr &  \le   &{1\over 2^{  m-n }\sqrt D}+   .... 
 .\end{eqnarray}
Let $\d$ be the arithmetical function defined by
\begin{eqnarray} \label{Delta}   \d(n)=\begin{cases}1\quad & {\rm if}\ n=1,\cr
0\quad & {\rm if}\ n \not=1. \end{cases} \end{eqnarray} Recall that 
\begin{eqnarray} \label{sp} \sum_{d|n}\m(d)=\d(n) . \end{eqnarray} 
 
 Consider first the sub-sum corresponding to $d=1$. We have  
\begin{eqnarray*}
 {1  \over  2^{  m-n }}
 \sum_{k=1\atop (k, D_{1/2})=1}^{m-n}   {{m-n}\choose {k}}
 &=&
 {1  \over   2^{  m-n }}
 \sum_{k=1}^{m-n}   {{m-n}\choose {k}}\d\big((k, D_{1/2})\big)
 \cr&=& 
 {1  \over D 2^{  m-n }}
 \sum_{k=1}^{m-n}   {{m-n}\choose {k}}\Big(\sum_{u|(k, D_{1/2} )}\m(u)\Big)
\cr&=& 
 \sum_{u|  D_{1/2} }\m(u)
 {1  \over   2^{  m-n }} \sum_{k=1\atop u|k}^{m-n}   {{m-n}\choose {k}}\cr&=& 
 \sum_{u|  D_{1/2} }\m(u)\,\big(\P\{u|B_{m-n}\}-{1  \over  2^{  m-n }}
\big)\end{eqnarray*}
Now let $(k, D_{1/2})=d>1$. Let $k=\k d$.
%, $D_{1/2}=\d d$
Then $(k,D_{1/2}/d)=1$, and so
\begin{eqnarray*}
& & {1  \over   2^{  m-n }}
 \sum_{k=1\atop (k, D_{1/2})=d}^{m-n}   {{m-n}\choose {k}}
\ = \ 
 {1  \over   2^{  m-n }}
 \sum_{1\le \k\le (m-n)/d\atop 
 (\k,D_{1/2}/d)=1}     {{m-n}\choose {\k d}}
\cr&=&
 {1  \over   2^{  m-n }}
 \sum_{1\le \k\le (m-n)/d}   {{m-n}\choose {\k d}}\d\big((\k, D_{1/2}/d)\big)
 \cr&=&{1  \over   2^{  m-n }}
 \sum_{1\le \k\le (m-n)/d}   {{m-n}\choose {\k d}}\Big(\sum_{u|(\k, D_{1/2}/d) }\m(u)\Big)
 \cr&=&\sum_{u|  D_{1/2}/d }\m(u)
 {1  \over   2^{  m-n }} \sum_{1\le \k\le (m-n)/d\atop u|\k}   {{m-n}\choose {\k d}}
    \cr &\buildrel{(\k = u\tau)}\over{=}& 
 \sum_{u|  D_{1/2} /d}\m(u)
 {1  \over   2^{  m-n }} \sum_{1\le ud\tau \le  m-n}     {{m-n}\choose {ud\tau}}
 \cr &=&
 \sum_{u|  (D_{1/2}/d ) }\m(u)
 {1  \over   2^{  m-n }} \sum_{1\le x \le  m-n\atop 
 ud|x}    {{m-n}\choose {x}}\cr&=& 
  \sum_{u|  (D_{1/2}/d )}\m(u)\,\big(\P\{ud|B_{m-n}\}-{1  \over  2^{  m-n }}
\big)
  \cr&=&
  \sum_{u|  (D_{1/2}/d )}\m(u) \P\{ud|B_{m-n}\}-{\d(D_{1/2}/d ) \over  2^{  m-n }}.\end{eqnarray*} 
Consequently,
 \begin{eqnarray*}
& &{1  \over  2^{  m-n }}\sum_{d| D_{1/2}} 2^{\o (d)}d
 \sum_{k=1\atop (k, D_{1/2})=d}^{m-n}   {{m-n}\choose {k}}  
\cr &=&
 \sum_{d| D_{1/2}} 2^{\o (d)}d
 \Big(\sum_{u|  (D_{1/2}/d )}\m(u)\,\big(\P\{ud|B_{m-n}\}-{\d(D_{1/2}/d )  \over  2^{  m-n }}\Big)
 \cr &=&
 \sum_{d| D_{1/2} } 2^{\o (d)}d
 \Big(\sum_{u|  (D_{1/2}/d )}\m(u)\, \P\{ud|B_{m-n}\}\Big)-{1  \over  2^{  m-n }} .
 \end{eqnarray*}
We observe   that $\d(D_{1/2}/d ) =0$ as long as $d< D_{1/2}$, and if $d=D_{1/2}$, then the corresponding summand   produces the contribution $2^{\o (D_{1/2})}D_{1/2}   \P\{ D_{1/2}|B_{m-n}\}-{1  \over  2^{  m-n }}$  (recalling that $\m(1)=1$), whence the last line. 
\vskip 3 pt
By reporting this in \eqref{323} we get
\begin{eqnarray}\label{323a} 
 {1\over D2^{ m-n }  } \sum_{k=0}^{m-n}  {{m-n}\choose {k}} \, \varrho_k(D) 
 & =  & {1  \over D 2^{  m-n }}\sum_{d| D_{1/2}} 2^{\o (d)}d
 \sum_{k=1\atop (k, D_{1/2})=d}^{m-n}   {{m-n}\choose {k}}  +{1\over 2^{  m-n }\sqrt D }
\cr & =  &  {1  \over D }\sum_{d| D_{1/2} } 2^{\o (d)}d
 \Big(\sum_{u|  (D_{1/2}/d )}\m(u)\, \P\{ud|B_{m-n}\}\Big)
 \cr & &  +{1\over 2^{  m-n }\sqrt D }-{1  \over  D2^{  m-n }} 
   .\end{eqnarray}
By inserting now this into  \eqref{4.11a}, we finally get
\begin{eqnarray}\label{4.11b}  \P \{ D|B_nB_m \} 
&\le &{1  \over D }\sum_{d| D_{1/2} } 2^{\o (d)}d
 \Big(\sum_{u|  (D_{1/2}/d )}\m(u)\, \P\{ud|B_{m-n}\}\Big)
\cr &  &   +{1\over 2^{  m-n }\sqrt D }-{1  \over D 2^{  m-n }}+C_\e
\left({     D  ^{1 + \e}\over   n }\right)^{1/2 }
\cr &\le &{1  \over D }\sum_{d| D_{1/2} } 2^{\o (d)}d
 \Big(\sum_{u|  (D_{1/2}/d )}\m(u)\, \P\{ud|B_{m-n}\}\Big)
\cr &  &   +{1\over 2^{  m-n }\sqrt D } +C_\e
\left({     D  ^{1 + \e}\over   n }\right)^{1/2 }. 
\end{eqnarray} 
This achieves the proof of Theorem \ref{bnmbound}.
\end{proof}

\end{document}
